\documentclass[10pt,a4paper]{article}
\usepackage[margin=1.8cm]{geometry}
\usepackage[T1]{fontenc}
\usepackage{lmodern}
\usepackage{booktabs}
\usepackage{tabularx}
\usepackage{array}
\usepackage{enumitem}
\usepackage[table]{xcolor}
\usepackage{titlesec}

\definecolor{accent}{HTML}{2B4FA3}
\definecolor{soft}{HTML}{EEF1F9}
\titleformat{\section}{\large\bfseries\color{accent}}{\thesection}{0.6em}{}
\titlespacing*{\section}{0pt}{10pt}{4pt}
\setlength{\parindent}{0pt}
\setlength{\parskip}{4pt}
\setlist{nosep, leftmargin=1.2em}
\newcolumntype{L}{>{\raggedright\arraybackslash}X}
\renewcommand{\arraystretch}{1.3}
\pagestyle{empty}

\begin{document}

{\LARGE\bfseries Emotion-Adaptive Generative UI}\\[2pt]
{\large Study queries and participant output}\\[-4pt]
{\color{accent}\rule{\linewidth}{0.8pt}}

\textbf{Setup.} A student uses \textbf{Arrow}, an AI chat-assistant app (sidebar, conversation thread, message box), and asks it
five fixed questions. Each answer was generated once by an LLM, reviewed by the team and frozen, so every participant
sees the same content. Every answer exists in \textbf{six looks}. A laptop webcam estimates facial emotion (our YOLO11n
model) and gaze (EyeTrax). When the student looks negative (confused or unhappy) for 4 seconds, the answer switches to
another look and asks \emph{``Keep this view / Go back''}. The student can also press \emph{``Show this differently''} at
any time. One session takes about 13--15 minutes.

\section{The five queries}
The queries are five different kinds of question on light, everyday topics (mostly computing), so the data covers
different kinds of generated UI and not just one kind of answer.

\begin{center}
\small
\rowcolors{2}{soft}{white}
\begin{tabularx}{\linewidth}{@{}c p{2.1cm} L p{1.9cm} L p{2.5cm}@{}}
\toprule
\textbf{\#} & \textbf{Topic} & \textbf{Query} & \textbf{Kind} & \textbf{What the answer contains} & \textbf{Quiz answer} \\
\midrule
Q1 & Machine learning & List the different machine learning models used for data analysis. \emph{(instructor's example)}
   & Comparison & model-choice tree, comparison table, model cards & k-means clustering \\
Q2 & Web \& networks & What happens when I type a website address into my browser?
   & Concept & flow diagram, steps, timing chart, ``ordering food'' analogy & Finds the IP address of the website's server \\
Q3 & Computer systems & Why does my computer slow down when I open too many apps?
   & Cause and effect & memory (RAM) chart, cause-effect flow, worked laptop example & To the storage drive (SSD or hard disk) \\
Q4 & Study skills & How should I prepare for my final exams in 2 weeks?
   & Plan & 14-day timeline, steps, evidence chart (testing beats re-reading) & Testing yourself with practice questions \\
Q5 & Grades & How do I calculate my GPA?
   & Calculation & formula, worked semester table, quality-points chart & 3.33 \\
\bottomrule
\end{tabularx}
\end{center}

\textbf{The six looks} (the same facts, presented differently):
\begin{center}
\small
\begin{tabularx}{\linewidth}{@{}p{2.6cm} L p{2.6cm} L@{}}
\toprule
\textbf{Look} & \textbf{Shows} & \textbf{Look} & \textbf{Shows} \\
\midrule
Paragraphs & a written explanation & Worked example & a concrete example with numbers \\
Visual & charts and diagrams & Summary & a short table or cheat sheet \\
Step by step & numbered steps & Simple & plain words and an everyday analogy \\
\bottomrule
\end{tabularx}
\end{center}

\textbf{Counterbalancing.} Query order follows a $5\times5$ Latin square across participants. The starting look is rotated,
so each look starts equally often.

\section{What happens in one query (about 2 minutes)}
\begin{enumerate}
  \item The question is pre-filled in the assistant box. The student presses \textbf{Ask} and the answer appears in its starting look.
  \item The student reads it. If the face is negative for 4\,s, the look switches automatically (at most 2 times per
        query, with a 15\,s gap) and the student chooses \emph{Keep} or \emph{Go back}. They can also switch it themselves.
  \item The student presses \textbf{I'm done}, then answers: satisfaction (1--7), which look helped most (if they saw
        more than one) and one quiz question.
\end{enumerate}

\section{Output for each participant}
Each participant gets one folder, \texttt{data/P01/}, with these files. All files share one clock.

\begin{center}
\small
\rowcolors{2}{soft}{white}
\begin{tabularx}{\linewidth}{@{}p{3.6cm} L@{}}
\toprule
\textbf{File} & \textbf{Contents} \\
\midrule
\texttt{interactions.csv} & \textbf{Main result: 5 rows, one per query.} Starting look, looks seen, final look,
   automatic and manual switches, kept/back counts, reading time and time per look, satisfaction, preferred look,
   quiz answer and correct (0/1), and emotion during the query (mean valence and arousal, dominant emotion,
   share of negative samples, emotion per look). \\
\texttt{adaptations.csv} & \textbf{One row per look switch} (training data for the future agent): trigger
   (auto/manual), emotion snapshot just \emph{before} the switch (valence, arousal, 8 emotion probabilities), from
   look $\rightarrow$ to look, response (keep / back / no answer), response time, and that query's
   satisfaction and quiz result. \\
\texttt{emotion.csv} & Facial emotion about 10 times a second: 8 class probabilities, valence and arousal. \\
\texttt{gaze.csv} & Where on the screen the student looked (webcam gaze). \\
\texttt{aoi.csv} & On-screen rectangle of every answer block (chart, table, text\ldots), used to tell what was looked at. \\
\texttt{looks.csv} & Every look change with its time and reason (initial, auto, manual, back). \\
\texttt{events.csv}, \texttt{mouse.csv} & Clicks, scrolling, button presses and mouse movement. \\
\texttt{questionnaires.csv} & Consent, age, programme, study confidence, all ratings and quiz answers, favourite
   look and comment. \\
\texttt{markers.csv} & Start and end time of every step and query. \\
\texttt{merged.csv} & Everything above on one 10\,Hz timeline (query, look, emotion, gaze block, mouse). \\
\texttt{session.json} & Session details: order, starting looks, settings, content version. \\
\bottomrule
\end{tabularx}
\end{center}

No video is saved, only numbers. \texttt{python run.py export} combines all participants into
\texttt{interactions\_all.csv}, \texttt{adaptations\_all.csv} and \texttt{dataset\_all.csv}.

\textbf{In total:} 9 participants $\times$ 5 queries = \textbf{45 interaction rows} in the first round (later about 125
students), plus one adaptation row for every switch. This tells us \emph{which look, under which emotion, helped the
student or not}, which is what the Generative-UI agent will be trained on.

\end{document}
